% GENERATED FILE. Edit canonical lessons, then run npm run generate:books.
% canonical-source-hash: fnv1a64:1e46cad4913151d6
% canonical-lessons: AR-C230-nabaha, AR-C230-atama, AR-C230-adafa, AR-C230-qassama, AR-C230-rakala

\chapter{To bark, To feed, To add, To divide, To kick}
\label{ch:ar-to-bark-to-feed-to-add-to-divide-to-kick}
\hlchaptermodality{230}

\begin{chapteropening}
\textbf{By the end of this chapter:} \emph{I can say to bark, to feed, to add, to divide and to kick.}

Complete the last lesson of chapter 230: I can say to bark, to feed, to add, to divide and to kick.
\end{chapteropening}

\section[\texorpdfstring{\ar{نبح} (nabaḥa)}{نبح (nabaḥa)}]{\ar{نبح} (nabaḥa) — to bark}

\label{lesson:AR-C230-nabaha}



\begin{quote}
Before the new one: say the Arabic for to knead, then the Arabic for to wipe.
\end{quote}

\subsection*{You'll want to know: \ar{نبح}}
\textbf{\ar{نبح}} — \emph{nabaḥa} — \textquotedblleft{}to bark\textquotedblright{}.

\begin{cousinweb}
\textbf{\ar{نبح}} (\emph{nabaḥa}) is built on the native Arabic root \emph{n-b-ḥ}, \textquotedblleft{}to bark\textquotedblright{} (Lane). \emph{Nabaḥa} is the plain verb, \textquotedblleft{}it barked\textquotedblright{}.
\end{cousinweb}

\begin{grammarlens}[title={What you've built}]
One more word for this chapter.
\end{grammarlens}

\subsection*{Your turn}
\begin{itemize}
\raggedright
  \item \emph{Say it:} \emph{nabaḥa}
  \item \emph{Say it:} \emph{nabaḥa}, once more
  \item \emph{Say it:} say \emph{nabaḥa}
  \item \emph{Recall:} say the Arabic for to knead, then the Arabic for to wipe, then say \emph{nabaḥa} again
\end{itemize}

\begin{tcolorbox}[breakable,colback=teal!4,colframe=teal!35!black,title={Before you move on}]
What does \ar{نبح} mean? (\textquotedblleft{}To bark\textquotedblright{}.) Say it once more.
\end{tcolorbox}
\section[\texorpdfstring{\ar{أطعم} (aṭʿama)}{أطعم (aṭʿama)}]{\ar{أطعم} (aṭʿama) — to feed}

\label{lesson:AR-C230-atama}



\begin{quote}
Before the new one: say the Arabic for to wipe, then the Arabic for to bark.
\end{quote}

\subsection*{You'll want to know: \ar{أطعم}}
\textbf{\ar{أطعم}} — \emph{aṭʿama} — \textquotedblleft{}to feed\textquotedblright{}.

\begin{cousinweb}
\textbf{\ar{أطعم}} (\emph{aṭʿama}) is built on the native Arabic root \emph{ṭ-ʿ-m}, \textquotedblleft{}to eat, to taste\textquotedblright{} (Lane). \emph{Aṭʿama} is Form IV, \textquotedblleft{}to make eat\textquotedblright{}, to feed.
\end{cousinweb}

\begin{grammarlens}[title={What you've built}]
One more word for this chapter.
\end{grammarlens}

\subsection*{Your turn}
\begin{itemize}
\raggedright
  \item \emph{Say it:} \emph{aṭʿama}
  \item \emph{Say it:} \emph{aṭʿama}, once more
  \item \emph{Say it:} say \emph{aṭʿama}
  \item \emph{Recall:} say the Arabic for to wipe, then the Arabic for to bark, then say \emph{aṭʿama} again
\end{itemize}

\begin{tcolorbox}[breakable,colback=teal!4,colframe=teal!35!black,title={Before you move on}]
What does \ar{أطعم} mean? (\textquotedblleft{}To feed\textquotedblright{}.) Say it once more.
\end{tcolorbox}
\section[\texorpdfstring{\ar{أضاف} (aḍāfa)}{أضاف (aḍāfa)}]{\ar{أضاف} (aḍāfa) — to add}

\label{lesson:AR-C230-adafa}



\begin{quote}
Before the new one: say the Arabic for to bark, then the Arabic for to feed.
\end{quote}

\subsection*{You'll want to know: \ar{أضاف}}
\textbf{\ar{أضاف}} — \emph{aḍāfa} — \textquotedblleft{}to add\textquotedblright{}.

\begin{cousinweb}
\textbf{\ar{أضاف}} (\emph{aḍāfa}) is built on the native Arabic root \emph{ḍ-y-f}, \textquotedblleft{}to lodge as a guest\textquotedblright{} (Lane). \emph{Aḍāfa} is Form IV, \textquotedblleft{}to attach, to add\textquotedblright{}.
\end{cousinweb}

\begin{grammarlens}[title={What you've built}]
One more word for this chapter.
\end{grammarlens}

\subsection*{Your turn}
\begin{itemize}
\raggedright
  \item \emph{Say it:} \emph{aḍāfa}
  \item \emph{Say it:} \emph{aḍāfa}, once more
  \item \emph{Say it:} say \emph{aḍāfa}
  \item \emph{Recall:} say the Arabic for to bark, then the Arabic for to feed, then say \emph{aḍāfa} again
\end{itemize}

\begin{tcolorbox}[breakable,colback=teal!4,colframe=teal!35!black,title={Before you move on}]
What does \ar{أضاف} mean? (\textquotedblleft{}To add\textquotedblright{}.) Say it once more.
\end{tcolorbox}
\section[\texorpdfstring{\ar{قسم} (qassama)}{قسم (qassama)}]{\ar{قسم} (qassama) — to divide, to share out}

\label{lesson:AR-C230-qassama}



\begin{quote}
Before the new one: say the Arabic for to feed, then the Arabic for to add.
\end{quote}

\subsection*{You'll want to know: \ar{قسم}}
\textbf{\ar{قسم}} — \emph{qassama} — \textquotedblleft{}to divide, to share out\textquotedblright{}.

\begin{cousinweb}
\textbf{\ar{قسم}} (\emph{qassama}) is built on the native Arabic root \emph{q-s-m}, \textquotedblleft{}to divide\textquotedblright{} (Lane). \emph{Qassama} is Form II, \textquotedblleft{}to divide up\textquotedblright{}.
\end{cousinweb}

\begin{grammarlens}[title={What you've built}]
One more word for this chapter.
\end{grammarlens}

\subsection*{Your turn}
\begin{itemize}
\raggedright
  \item \emph{Say it:} \emph{qassama}
  \item \emph{Say it:} \emph{qassama}, once more
  \item \emph{Say it:} say \emph{qassama}
  \item \emph{Recall:} say the Arabic for to feed, then the Arabic for to add, then say \emph{qassama} again
\end{itemize}

\begin{tcolorbox}[breakable,colback=teal!4,colframe=teal!35!black,title={Before you move on}]
What does \ar{قسم} mean? (\textquotedblleft{}To divide, to share out\textquotedblright{}.) Say it once more.
\end{tcolorbox}
\section[\texorpdfstring{\ar{ركل} (rakala)}{ركل (rakala)}]{\ar{ركل} (rakala) — to kick}

\label{lesson:AR-C230-rakala}



\begin{quote}
Before the new one: say the Arabic for to add, then the Arabic for to divide.
\end{quote}

\subsection*{You'll want to know: \ar{ركل}}
\textbf{\ar{ركل}} — \emph{rakala} — \textquotedblleft{}to kick\textquotedblright{}.

\begin{cousinweb}
\textbf{\ar{ركل}} (\emph{rakala}) is built on the native Arabic root \emph{r-k-l}, \textquotedblleft{}to kick\textquotedblright{} (Lane). \emph{Rakala} is the plain verb, \textquotedblleft{}he kicked\textquotedblright{}.
\end{cousinweb}

\begin{grammarlens}[title={What you've built}]
One more word for this chapter.
\end{grammarlens}

\subsection*{Your turn}
\begin{itemize}
\raggedright
  \item \emph{Say it:} \emph{rakala}
  \item \emph{Say it:} \emph{rakala}, once more
  \item \emph{Say it:} say \emph{rakala}
  \item \emph{Recall:} say the Arabic for to add, then the Arabic for to divide, then say \emph{rakala} again
\end{itemize}

\begin{tcolorbox}[breakable,colback=teal!4,colframe=teal!35!black,title={Before you move on}]
What does \ar{ركل} mean? (\textquotedblleft{}To kick\textquotedblright{}.) Say it once more.
\end{tcolorbox}
